\subsection{问题一：有界误差下的交会定位}

记 $u(\alpha)=(\cos\alpha,\sin\alpha)^\mathsf{T}$，二维叉积为
$a\times b=a_xb_y-a_yb_x$。在检测点 $S_i$ 测得示向度
$\theta_i$ 后，真实方位角落在
$[\theta_i-\varepsilon,\theta_i+\varepsilon]$。因此候选点 $p$ 必须满足
\begin{equation}
\begin{split}
u(\theta_i-\varepsilon)\times(p-S_i)&\ge 0,\\
u(\theta_i+\varepsilon)\times(p-S_i)&\le 0.
\end{split}
\label{eq:wedge}
\end{equation}
每条观测给出两个线性半平面。将全部约束写成 $Ap\le b$，得到
\begin{equation}
\Omega=\{p\in\mathbb{R}^2:Ap\le b\}.
\label{eq:region}
\end{equation}
先用线性规划判定式\eqref{eq:region}是否可行，再分别沿 $\pm x$、$\pm y$ 方向求线性目标；任一方向无界即判定定位区域无界。对于有界区域，枚举两条非平行边界线的交点并检查全部约束，去重后按凸包顺序得到顶点集
$V=\{v_1,\ldots,v_m\}$。凸函数 $\|p-q\|_2$ 在紧凸多边形乘积
$\Omega\times\Omega$ 上的最大值可在顶点对取得，故
\begin{equation}
D(\Omega)=\max_{1\le i<j\le m}\|v_i-v_j\|_2.
\label{eq:diameter}
\end{equation}
直接枚举复杂度为 $O(m^2)$，规模较大时可在凸包上用旋转卡壳降至
$O(m)$\cite{orourke1998}。空集、无界集与退化为线段或单点的情况分别返回状态，避免用人工大矩形产生虚假有限直径。

仅知道直径 $D$ 不能保证存在半径 $D/2$ 的覆盖圆。我们进一步枚举由一个、两个或三个顶点支撑的圆，选取覆盖全部顶点的最小半径，得到最小包围圆 $(c_*,r_*)$；该问题也可用随机增量算法在线性期望时间内求解\cite{welzl1991}。只有当 $r_*\le20\unit{m}$ 时，移动到 $c_*$ 才能在清除半径内覆盖整个定位区域。

设观测数为 $n$，半平面数为 $2n$。本文实现枚举边界线对并逐条检查约束，最坏时间复杂度为 $O(n^3)$；得到 $m$ 个凸包顶点后，直径基线为 $O(m^2)$。最小包围圆的确定性核验枚举至多三个支撑点，并用全部顶点复查，最坏为 $O(m^4)$。这些算法用于观测次数很少的单源定位，计算耗时远小于一次 $5\unit{s}$ 的测量动作，因而优先采用易于审计的确定性写法。实现中对平行线、重复顶点和半平面边界统一使用 $10^{-8}$ 量级容差，输出后再用全部原始约束复验，容差只处理浮点舍入，不放宽 $20\unit{m}$ 的物理清除条件。

\subsection{问题二：保证接收的第二检测点}

设首次检测点为 $S_1$，返回示向度 $\theta$。全向源的可能位置保守写为
\begin{equation}
\Omega_1=\{S_1+\rho u(\alpha):0\le\rho\le R,
\ |\alpha-\theta|\le\varepsilon\},\qquad R=1500\unit{m}.
\label{eq:first-sector}
\end{equation}
返回 \texttt{direction} 实际还说明 $\rho>5\unit{m}$，但在保证性设计中保留扇形顶点只会扩大可行域。记远端两点
$E_\pm=S_1+R u(\theta\pm\varepsilon)$。要在未知实际接收半径下保证第二次收到信号，第二点 $S_2$ 到任一 $G\in\Omega_1$ 的距离必须不超过下界 $r_0=1000\unit{m}$。

对固定 $S_2$，沿每条半径的距离平方是关于 $\rho$ 的凸二次函数，其最大值在 $\rho=0$ 或 $R$ 取得；在远端圆弧上，最大距离在两个端点之一取得。因此保证接收的候选区域恰为
\begin{equation}
\mathcal C=B(S_1,r_0)\cap B(E_-,r_0)\cap B(E_+,r_0).
\label{eq:candidate}
\end{equation}
这一定义同时给出可计算判据：候选点到三个圆心的最大距离不超过
$1000\unit{m}$。

为避免两条示向线近似平行，在候选域上下两半分别取最大的横向基线。令
$v(\theta)=(-\sin\theta,\cos\theta)^\mathsf{T}$，并记
\begin{equation}
q=\sqrt{r_0^2-(R/2)^2}.
\end{equation}
取首测点圆与相反侧远端点圆的上、下交点，得到对称推荐式
\begin{equation}
S_2^\pm=S_1+a\,u(\theta)\pm d\,v(\theta),
\quad
a=\frac R2\cos\varepsilon+q\sin\varepsilon,
\quad
d=q\cos\varepsilon-\frac R2\sin\varepsilon.
\label{eq:second-point}
\end{equation}
式\eqref{eq:second-point}可由两个等半径圆的交点直接得到。在局部坐标中，
$S_1=(0,0)$，相反侧远端点为
$E_-=(R\cos\varepsilon,-R\sin\varepsilon)$，两圆心中点到交点的垂距为 $q$；把垂向分量旋转回 $(u,v)$ 基即可得到 $a,d$。因此推荐点到 $S_1$ 和一个远端点的距离恰为 $r_0$，再检查它到另一个远端点的距离不超过 $r_0$，即可判定属于三圆盘交集。

根式为实数要求 $R\le2r_0$，这是存在单点保证接收的必要条件：当
$R>2r_0$ 时，$B(S_1,r_0)$ 与任一远端圆盘已经不相交。即便满足该必要条件，仍须检查三个圆盘是否有公共交集以及 $d>0$；当 $d$ 接近零时，第二点虽然可能保证接收，但交会基线退化，方位误差会被显著放大。官方参数下 $q=661.438\unit{m}$、$d=648.248\unit{m}>0$，故闭式点同时满足接收和非共线要求。

若第二点返回示向度，则按式\eqref{eq:wedge}把新楔形与
$\Omega_1$ 相交；若返回 \texttt{near}，则已进入 $5\unit{m}$ 邻域，可直接执行光学定位清除。对于其他参数，也可在 $\mathcal C$ 内离散搜索，使所有可能源位置下的最小交会角最大，或使更新后定位区域的最坏直径最小。

\subsection{问题三：全向源的保证覆盖搜索}

取边长 $h=600\unit{m}$ 的正方形网格。保留所有与半径
$1800\unit{m}$ 源域相交方格的四个顶点，共得到 61 个去重覆盖点。任一源位于某个被保留方格中，它到该方格至少一个顶点的距离不超过
\begin{equation}
\sqrt{2}h=848.53\unit{m}<1000\unit{m}.
\label{eq:grid-cover}
\end{equation}
全向源在该点必能被接收。机器狗在每个覆盖点扫描全部尚未清除的频道；定位清除产生绕行后，从当前位置选择最近的未完成覆盖点继续，覆盖点集合本身不改变。

覆盖点集不是在圆内等间距抽样，而是先保留与源圆域相交的闭方格，再取这些方格的顶点并去重。形式化地，令
\begin{equation}
\mathcal Q=\{[ih,(i+1)h]\times[jh,(j+1)h]:Q_{ij}\cap B(0,1800)\ne\varnothing\},
\qquad
\mathcal P=\bigcup_{Q\in\mathcal Q}\operatorname{vert}(Q).
\label{eq:coverage-set}
\end{equation}
代码只需在 $i,j\in\{-4,\ldots,3\}$ 的有限范围内检查方格到原点的最短距离，最终 $|\mathcal P|=61$。边界方格的部分顶点位于半径 $1800\unit{m}$ 圆外，但不能删除：它们可能是边界源到最近的接收点，也是问题四中落入辐射半平面的唯一见证顶点。

某频道首次返回正示向后，以半径 $1500\unit{m}$ 的示向扇形与源圆域外逼近作为初始定位区域。后续在最小包围圆中心两侧自适应插入检测点，每次只用正示向更新区域。若最小包围圆半径满足
$r_*\le19.5\unit{m}$，则在圆心执行清除；$0.5\unit{m}$ 为数值裕量。最多八次自适应测量后仍未达到阈值时，以边长 $20\unit{m}$ 网格中心覆盖剩余多边形并按最近点顺序逐点清除，任一可行点到某网格中心的距离不超过
$10\sqrt2\unit{m}<20\unit{m}$，故兜底过程有限且完备。

记移动总距离为 $L$。检测和频道切换次数分别为 $N_m,N_s$，成功和失败清除次数分别为
$N_c,N_f$。虚拟总时间按协议复算为
\begin{equation}
T=\frac L5+5N_m+N_s+5N_c+3N_f.
\label{eq:time}
\end{equation}
当 61 个点全部完成时，覆盖证明给出无遗漏；当已清除 16 个源时，利用题面数量上界可立即结束。除此之外不以“已找到 10 个”作为停止条件。

将覆盖点访问、单源定位和清除看作离散状态之间的规划过程\cite{lavalle2006}，可写成如下有限步骤：
\begin{enumerate}
  \item 从当前位置选择最近的未完成覆盖点，依次扫描尚未清除的频道；
  \item 频道返回正示向时，建立“示向楔形与源圆域”的初始交集；返回无信号时只继续扫描，不删除定向源候选位置；
  \item 在当前可行域上自适应选择下一检测点，用新的正示向相交更新；
  \item 若 $r_*\le19.5\unit{m}$，移动到最小包围圆圆心清除；八次后仍不满足时，执行 $20\unit{m}$ 网格兜底；
  \item 清除成功后继续最近未完成覆盖点；清除数达到 16 或全部 61 个覆盖点完成时退出并保存日志。
\end{enumerate}

图\ref{fig:algorithm-flow}展示外层覆盖与内层定位之间的控制关系。流程中的每个循环要么完成一个覆盖点，要么清除一个频道对应的源；覆盖点和频道数均有限，单源自适应测量又有八次上限，因此算法必在有限步内终止。

更具体地说，若没有提前清除 16 个源，外层至多弹出 61 个覆盖点；每个覆盖点至多访问 20 个尚未清除频道。每个首次发现的频道至多进行八次自适应测量，随后进入对当前有界多边形的有限网格覆盖。网格边长为 $20\unit{m}$，多边形又被半径 $1800\unit{m}$ 源域的外逼近限制，故兜底清除点数也有限。该论证给出终止性和发现完备性，但不声称路线长度或虚拟时间达到全局最优。

\begin{figure}[htbp]
  \centering
  \includegraphics[width=0.98\textwidth]{q34-algorithm-flow.pdf}
  \caption{问题三、四共用的覆盖搜索、定位与清除流程}
  \label{fig:algorithm-flow}
\end{figure}

\subsection{问题四：混合源的发现完备性}

定向源只在通过源点、法向由发射方向确定的闭半平面内辐射。对包含源
$G$ 的覆盖方格，$G$ 位于四顶点的凸包内。任一通过 $G$ 的闭半平面至少包含一个方格顶点，否则四顶点都严格位于另一侧，其凸包也不可能包含
$G$。结合式\eqref{eq:grid-cover}，至少一个顶点既在辐射半平面内又在
$1000\unit{m}$ 距离内，因此同一 600 米覆盖点集仍能保证发现定向源。

定向源的 \texttt{no\_signal} 可能由距离过远或位于背面导致，不能据此删去位置区域；算法只使用正示向相交定位。清除成功只由清除点与源的距离决定，与定向方向无关，故发现后的定位、最小包围圆清除和有限网格兜底与问题三共用。

图\ref{fig:coverage-proof}左侧给出源圆域与 61 个覆盖点，右侧放大一个边长 $600\unit{m}$ 的方格。任意辐射半平面至少包含一个着色顶点，该顶点到方格内源点的距离不超过对角线 $848.53\unit{m}$，因而同时满足方向和距离条件。

\begin{figure}[htbp]
  \centering
  \includegraphics[width=0.98\textwidth]{q34-coverage-directional-proof.pdf}
  \caption{600 米方格覆盖及定向源可接收见证点}
  \label{fig:coverage-proof}
\end{figure}

通信层依次执行 \texttt{enter}、\texttt{measure/clear} 和 \texttt{exit}。网络状态不确定时，重试严格复用同一请求内容与请求编号。
